\documentclass[11pt,a4paper]{article}
\usepackage[T1]{fontenc}
\usepackage{lmodern,microtype,amsmath,amssymb,booktabs,tabularx,array}
\usepackage[margin=22mm]{geometry}
\usepackage{xcolor,hyperref}
\definecolor{navy}{HTML}{193B54}
\hypersetup{colorlinks=true,linkcolor=navy,urlcolor=navy,
pdftitle={Supervised spectral Tucker: matched factor fine-tuning experiment}}
\setlength{\parindent}{0pt}\setlength{\parskip}{6pt}
\newcolumntype{Y}{>{\raggedright\arraybackslash}X}
\newcommand{\code}[1]{{\small\nolinkurl{#1}}}
\begin{document}
{\LARGE\bfseries\color{navy} Supervised spectral Tucker}\par
{\large A matched test of classification-driven factor learning}\par
Successor to the tensor temporal screen; designed 4 October 2026.

\section{Hypothesis and controlled intervention}
The previous single-seed screen found that frozen spectral Tucker retained some
classification performance under severe electrode loss, but its full-input
validation balanced accuracy was 45.47\%, below the spectral dense control's
50.97\% and the EEGNet--Transformer reference's 63.71\%. These results motivate
one bounded hypothesis: updating Tucker factors for classification, while
preserving observed-entry ridge inference, may improve discriminative features.

This study compares two otherwise identical spectral Tucker models. Both start
with the same reconstruction-fitted factors and the same randomly initialized
temporal classifier. The supervised arm allows cross-entropy gradients to update
the factors; the frozen arm retains them as buffers. We use the same direct
ridge-solve implementation in both arms, so a change of numerical solver is
not confounded with supervision. This is factor learning during classifier
training, not fine-tuning an already selected classifier checkpoint.

\begin{tabularx}{\linewidth}{@{}lYY@{}}\toprule
Property & Frozen control & Supervised factors \\ \midrule
Spectral features and normalization & Fixed; training-calibrated & Identical \\
Initial Tucker factors & Training reconstruction fit & Identical fit and seed \\
Initial head and data order & Paired & Paired \\
Core inference & Observed-entry ridge solve & Same solve, differentiable \\
Factors during classification & Frozen buffers & Trainable $U,V$ \\
Factor constraint & Initial unit columns retained & Unit columns after updates \\
Classification loss & Cross-entropy & Same cross-entropy \\
Training availability & Full 22 channels & Full 22 channels \\ \bottomrule
\end{tabularx}

All comparisons in this report concern these two new fits. The previous
EEGNet--Transformer score is historical context, not a newly matched arm in
this experiment. A gain would support task-informed factor adaptation under
the tested settings, rather than proving a universal advantage of tensor models.

\section{Pipeline and differentiable core}
BCI IV 2a supplies four motor-imagery classes and 22 EEG electrodes at 250 Hz.
We retain the verified participant-specific T-session splits and the previous
artifact exclusions, cue alignment, 2--30 Hz window-local filtering and
training-only channel normalization. Each trial has four 250-sample windows.
Four overlapping 100-sample patches per window (stride 50) give 16 ordered time
positions without crossing window or availability boundaries.

A periodic Hann window and FFT yield 12 log-power features at 2.5--30 Hz.
The feature mean and population standard deviation are fitted on training trials
only, separately for each channel/frequency. With patch index $t$, standardized
features have shape $F[B,22,16,12]$. Boolean availability $M[B,22,4]$ expands
to the 16 patch positions. Hidden values are selected away before arithmetic,
and every original window must retain at least one observed channel.

Reconstruction fitting initializes $U\in\mathbb R^{22\times4}$ and
$V\in\mathbb R^{12\times4}$ for 30 epochs, using training data only. For a patch,
the inferred core solves
\[
 G^*=\arg\min_G\frac{\|F_O-U_OGV^\top\|_F^2}{12|O|}
             +\lambda\|G\|_F^2,\qquad \lambda=10^{-3}.
\]
Let $A=U_O^\top U_O$, $B=V^\top V$, $C=U_O^\top F_OV$ and
$\rho=12|O|\lambda$. Then
\[
 AG B+\rho G=C.
\]
Using row-major vectorization, the $16\times16$ system is
\[
 H_{(r,s),(r',s')}=A_{rr'}B_{ss'}+
       \rho\,\delta_{rr'}\delta_{ss'},\qquad
 g=H^{-1}c.
\]
Implementation uses \code{torch.linalg.solve} in float64, without explicitly
forming an inverse. Positive ridge makes $H$ positive definite even when fewer
than four electrodes are observed or factor columns correlate. Differentiating
the solve propagates
\[
 dg=H^{-1}(dc-dH\,g)
\]
to both factors. There is no detached core and no differentiation through an
eigendecomposition. The frozen and supervised forward functions agree at matched
initialization; the frozen result is also checked against the original solver.

Each $4\times4$ core is flattened to 16 values and projected to 32 features.
Fixed sinusoidal positions and projected availability bits enter a one-layer
temporal Transformer: four heads, feed-forward width 64, pre-LayerNorm, GELU
and dropout 0.1. Final normalization, mean pooling and 88 original mask flags
feed the four-class linear output. Spectral normalization remains fixed in both
arms; the learned factors are the only added trainable quantities.

Unit column norms are maintained by projection after supervised optimizer
updates. This is a norm constraint, not orthogonality. Reconstruction fitting
remains an initialization stage; classifier training adds no reconstruction loss.

\section{Training, evaluation and decision rule}
The loss is unweighted four-class cross-entropy. Both arms use AdamW with
learning rate $10^{-3}$, weight decay $10^{-4}$, batch size 32 and gradient norm
clipping at 1. The same learning rate applies to the head and supervised factors.
Ten warm-up epochs precede cosine decay. Fits run for at most 250 epochs;
early stopping is permitted after 75 epochs with patience 50. Full-input
validation balanced accuracy selects the checkpoint, breaking ties by lower
log loss and then earlier epoch.

There are nine participants and three paired split/training seeds (0, 1, 2):
27 paired tasks and 54 fits. Each seed changes both the historical split and
training randomness; these are not three restarts of one fixed split. No test
predictions or scores enter fitting, checkpoint selection or this analysis.

Validation includes full availability and static/dynamic random retention of
16 or 6 electrodes. Each degraded condition has five paired repeats. Aggregation
averages repeats, then seeds within participant, then participants equally.
The primary contrast is supervised minus frozen full-input balanced accuracy.
The exploratory screen requires at least two percentage points mean improvement
and improvement in at least six of nine participant means. Participant bootstrap
intervals use 2,000 resamples and seed 20261004. Seed-level findings, log loss,
train--validation gaps and degraded conditions are secondary evidence.

Three seeds improve the assessment of split/training variation, but do not make
validation independent of earlier research. An independent session or dataset
would be needed for confirmation. Condition-specific intervals are exploratory
and do not correct for multiple comparisons. All arms use full-input training;
mixed-availability training remains a separate hypothesis.

Source/configuration/package/input hashes identify the new study. Calibration
and initial-state hashes check pair matching, selected checkpoints are reloaded,
and artifacts are verified before reuse. The previous studies' source files,
configurations and scientific artifacts are preserved.

\section{Results and verdict}
All 54 real-data fits completed (nine participants, three paired split/training
seeds, two arms). The following values average seeds within participant and
participants equally. Accuracy values are percentages; differences are percentage
points. Log loss is unitless and lower is better.

\begin{center}\begin{tabular}{@{}lrrrr@{}}\toprule
Model & Train BA & Validation BA & Gap & Val. log loss \\ \midrule
Frozen factors & 67.44 & 45.71 & 21.73 & 1.593 \\
Supervised factors & 72.86 & 47.47 & 25.40 & 1.570 \\
\bottomrule\end{tabular}\end{center}

\begin{center}\begin{tabular}{@{}lrrr@{}}\toprule
Participant & Frozen BA & Supervised BA & Difference \\ \midrule
A01 & 39.74 & 45.70 & +5.95 \\
A02 & 38.10 & 41.80 & +3.71 \\
A03 & 59.71 & 63.42 & +3.71 \\
A04 & 42.36 & 41.93 & -0.43 \\
A05 & 40.38 & 41.03 & +0.64 \\
A06 & 42.42 & 42.58 & +0.15 \\
A07 & 41.85 & 43.04 & +1.19 \\
A08 & 48.53 & 49.08 & +0.55 \\
A09 & 58.32 & 58.62 & +0.31 \\
\bottomrule\end{tabular}\end{center}

\begin{center}\begin{tabular}{@{}lrrr@{}}\toprule
Availability & Frozen BA & Supervised BA & Difference \\ \midrule
Full 22 & 45.71 & 47.47 & +1.75 \\
Static 16 & 37.94 & 40.03 & +2.08 \\
Dynamic 16 & 37.61 & 40.47 & +2.86 \\
Static 6 & 35.42 & 35.38 & -0.04 \\
Dynamic 6 & 35.52 & 36.37 & +0.85 \\
\bottomrule\end{tabular}\end{center}

\textbf{Primary finding.} Updating the factors improved full-input balanced
accuracy by 1.75 percentage points, with an exploratory participant-bootstrap
95\% interval of $[0.56,3.18]$ points. Eight of nine participant averages improved;
A04 declined by 0.43 points. However, the prespecified two-point mean-gain
threshold was not met. The appropriate verdict is a small, fairly consistent
benefit from supervised factors, with a failed practical screening criterion.
It is not a null result, nor evidence of a strong classification architecture.

Split/training-seed mean gains were $+1.80$, $+2.58$ and $+0.87$ points for
seeds 0, 1 and 2. Seventeen of the 27 individual participant/seed pairs improved,
eight declined and two tied (absolute tolerance $10^{-12}$ in balanced accuracy).
Training balanced accuracy increased by 5.42 points while validation increased
by 1.75, widening the mean train--validation gap from 21.73 to 25.40 points.
This is consistent with additional fitting capacity exceeding the gain in
generalization; it does not uniquely identify the cause. Full-input log loss
changed by $-0.023$ with interval $[-0.243,0.182]$, providing no clear improvement
in probabilistic prediction. Both mean full-input log losses exceed the uniform
four-class predictor's $\log 4\simeq1.386$, despite above-chance balanced accuracy.
This indicates a weakness in confidence quality. All 54 fits stopped before
the 250-epoch limit (75--121 epochs); extending the epoch budget is not supported
by these histories.

\textbf{Electrode-loss findings.} At 16 retained channels, gains were
$+2.08$ points for static masks (interval $[0.81,3.10]$) and $+2.86$ for dynamic
masks ($[1.08,4.47]$), each positive in eight participants. At six channels,
static masks changed by $-0.04$ points ($[-1.28,1.45]$) and dynamic masks by
$+0.85$ ($[-0.89,2.72]$). Thus this intervention helps moderate-loss accuracy,
but does not establish a severe-loss benefit. Mean log loss at six channels
increased by 0.068 (static) and 0.027 (dynamic), with both intervals crossing
zero. All these condition-specific findings are exploratory.

\textbf{Research decision.} Keep differentiable factor learning as a modestly
supported component, but do not promote this model as the accuracy baseline.
The 47.47\% validation balanced accuracy remains below the earlier spatial
Transformer's 63.71\% single-seed screen and 64.87\% three-seed study; those are
historical context, not matched contrasts in this study. The result supports
classification-informed factor adaptation at these fixed spectral features,
ranks and head, but cannot show that task adaptation alone repairs the
information bottleneck. It also does not test learned raw temporal features.

A useful next bounded comparison would train both frozen and supervised arms
with the same mixed-availability schedule, retaining this full-training study
as a control. That would test whether exposure to missing channels improves
the six-channel outcome. For the separate full-input accuracy objective,
test a simpler or more strongly regularized classifier on the same cores before
increasing factor ranks or Transformer depth. The larger train--validation gap
motivates this test but does not prove that classifier capacity is the cause.
Each change
requires its own declared identity and controls; neither experiment was run here.

\textbf{Verification and provenance.} The independent audit replays all 54
selected checkpoints on training and validation data, reconstructs all 1,134
availability evaluations, recalibrates each task's training-only initial factors,
and independently checks aggregate means and bootstrap intervals. Selected
supervised factors changed in all 27 fits, while frozen factors were unchanged;
paired initial full states and calibration histories match exactly. The software
checks comprise 25 new model/runner tests and 17 tensor temporal regression tests.
The numerical review checks the direct system against an explicit observed-entry
design and verifies gradients. Historical source/configuration/workflow/result
preservation checks cover 1,122 files.

All 85 Python source hashes, 64 input hashes, configuration and package versions
are recorded in the manifest. The study identity is
\begin{center}{\scriptsize\ttfamily
479db2ec18e7f1e1b5d1e5255c5a200c3df7d59b7d8a73936b4415252af8ed57}
\end{center}
The machine-readable results and independent audit are under
\code{results/supervised-tucker-v1/report/}. Fits used one CPU thread each;
cohort execution used up to four independent processes. Elapsed times are
execution records, not controlled speed benchmarks. Checkpoint selection and
reporting used validation only, so an independent session or dataset remains
necessary before a confirmatory accuracy claim.


\section{Reproduction}
Declaration: \code{configs/supervised-tucker.json}.

Implementation: \code{inm/supervised_tucker/}.

Launcher: \code{scripts/run_supervised_tucker_experiments.sh}.

Use separate \code{--pilot}, \code{--cohort}, \code{--summarize} and
\code{--dry-run} modes.
Outputs belong to \code{results/supervised-tucker-v1/}.

Build this report with \verb|make -C docs/concept supervised_tucker.pdf|.
\end{document}
